\section[Cote 100 : un point de dimension un hors d'une famille dénombrable de diviseurs]{Cote 100 : un point de dimension un hors d'une famille dénombrable de diviseurs\protect\verifie{Folder100}}\label{sec:100}

La cote 100, \emph{Bonne réduction des variétés abéliennes via bonne réduction de Barsotti--Tate $T_p(A)$ (base quelconque)}, cinquante-deux pages que l'inventaire date \q{[à partir de 1967-1968]}, mêle des notes manuscrites et la copie annotée d'un tapuscrit. Pour algébriser un schéma abélien formel, les pages~30 à~38 ont besoin d'un lemme d'algèbre commutative : dans le spectre d'un anneau local noethérien complet intègre, une famille dénombrable de fermés rares laisse échapper un point $x$ avec $\dim \overline{\{x\}} = 1$ (corollaire de la page~38, \lot{100}{2}{38}). La page le démontre en passant à $\Theta = V[[T_1, \ldots, T_n]]$, $V$ anneau de valuation discrète complet. Un $V$-homomorphisme $\varphi : \Theta \to V$ est déterminé par les images des $T_i$, qui sont des éléments arbitraires de l'idéal maximal $\mathfrak m$ ; son noyau est un point de dimension un. Chaque fermé rare, contenu dans un $V(f)$ avec $f \neq 0$, exclut les $\varphi$ tels que $f(\varphi(T_1), \ldots, \varphi(T_n)) = 0$, partie fermée rare de l'espace métrique complet $\mathfrak m^n$, et le théorème de Baire conclut.

On démontre ici le cas $n = 1$, celui de $\Theta = V[[T]]$, anneau de dimension~$2$. Les homomorphismes $\Theta \to V$ sont alors les $\varphi_a : T \mapsto a$, $a \in \mathfrak m$, de noyau l'idéal $(T - a)$. Dans ce cas l'argument de Baire peut être remplacé par un dénombrement : chaque série $f \neq 0$ n'exclut qu'un nombre fini de valeurs de $a$, tandis que $\mathfrak m$ n'est pas dénombrable.

\noindent\emph{Conventions.} $V$ est un anneau de valuation discrète, complet pour la topologie $\mathfrak m$-adique, d'uniformisante $\varpi$ ; $V[[T]]$ est l'anneau des séries formelles. Dire qu'un point de dimension un, le noyau de $\varphi_a$, est hors de $V(f)$ revient à dire que $f$ n'appartient pas à l'idéal $(T - a)$.

\begin{theoreme}[énoncé de l'édition, cas $n = 1$ du corollaire de la page 38\piste{100-baire-dimension-one-points}]\label{thm:100-points}
Soit $V$ un anneau de valuation discrète complet. Pour $a \in \mathfrak m$, l'anneau $V[[T]]/(T - a)$ est isomorphe à $V$. Pour toute famille dénombrable $(f_i)$ de séries non nulles de $V[[T]]$, il existe $a \in \mathfrak m$ tel qu'aucune $f_i$ n'appartienne à $(T - a)$.
\end{theoreme}

\begin{proof}
Le polynôme $T - a$ est unitaire et son terme constant $-a$ est dans $\mathfrak m$ : il est distingué au sens de Weierstrass. La division de Weierstrass identifie alors $V[[T]]/(T - a)$ à $V[T]/(T - a)$, qui est $V$ par $P \mapsto P(a)$, et un polynôme $P$ est dans l'idéal $(T - a)$ de $V[[T]]$ si et seulement si $P(a) = 0$.

Soit $f \neq 0$. Le plus petit $v$ tel que $\varpi^{v+1}$ ne divise pas tous les coefficients de $f$ existe, et $f = \varpi^v g$ où $g$ a un coefficient inversible, c'est-à-dire que l'image de $g$ dans $k[[T]]$, $k$ le corps résiduel, n'est pas nulle. Le théorème de préparation de Weierstrass écrit $g = P u$, avec $P$ polynôme distingué et $u$ inversible. Donc $f = Q u$, où $Q = \varpi^v P$ est un polynôme non nul, et $f \in (T - a)$ si et seulement si $Q(a) = 0$ : seules les racines de $Q$, en nombre fini, sont exclues.

Il reste à voir que $\mathfrak m$ n'est pas dénombrable ; comme $a \mapsto \varpi a$ est une injection de $V$ dans $\mathfrak m$, il suffit que $V$ ne le soit pas. À toute suite $\varepsilon \in \{0, 1\}^{\mathbf N}$, la complétude associe une limite $L_\varepsilon$ des sommes partielles $\sum_{i < n} \varepsilon_i \varpi^i$, congrue à la $n$-ième modulo $\varpi^n$. Si $\varepsilon$ et $\varepsilon'$ diffèrent pour la première fois en $j$, les sommes partielles d'indice $j + 1$ diffèrent de $\pm\varpi^j$, et $L_\varepsilon = L_{\varepsilon'}$ forcerait $\varpi^{j+1} \mid \varpi^j$, ce qui est impossible. L'ensemble des $a \in \mathfrak m$ exclus par l'une des $f_i$ est une réunion dénombrable d'ensembles finis, donc dénombrable, et il ne peut pas contenir~$\mathfrak m$.
\end{proof}

Le dénombrement ne remplace Baire qu'en une variable. Pour $n \geq 2$, une série non nulle exclut en général une infinité non dénombrable de points de $\mathfrak m^n$, par exemple tout l'hyperplan $T_1 = 0$ pour $f = T_1$, et il faut l'argument de la page : une série non nulle ne s'annule identiquement sur aucun ouvert de $\mathfrak m^n$. Le théorème ne suppose rien sur le corps résiduel, qui peut être fini, comme pour $V = \mathbf Z_p$ : la non-dénombrabilité de $V$ vient de la complétude seule.

\noindent Vérifié tel quel (\texttt{quotEquiv}, \texttt{exists\_poly\_mul\_unit}, \texttt{finite\_bad}, \texttt{uncountable}, \texttt{exists\_avoid}), avec le théorème de préparation de Weierstrass et l'isomorphisme $V[T]/(g) \simeq V[[T]]/(g)$ de mathlib pour un polynôme distingué~$g$.

Ne sont pas démontrés ici : le cas de plusieurs variables, où Baire est nécessaire ; la réduction d'un anneau local complet intègre à une extension finie de $V[[T_1, \ldots, T_n]]$ (théorème de structure de Cohen) et la remontée des points de dimension un le long de cette extension ; la proposition de la page~36, qui s'en déduit. L'énoncé général reste ouvert, comme les autres pistes de la cote :

\begin{encadre}
\begin{theoreme*}[Grothendieck, pages 2, 3 et 20, pleine fidélité et composantes connexes\piste{100-full-faithfulness-connected-components}]\nonverifie
Soient $S$ un schéma localement noethérien et $U \subset S$ un ouvert tel que toutes les caractéristiques résiduelles de $S - U$ soient égales à $p$. Le foncteur $A \mapsto (A_U, A[p^\infty], \mathrm{id})$ des schémas abéliens sur $S$ vers les triplets est pleinement fidèle dès que $U$ rencontre chaque composante connexe de $S$, sans hypothèse de normalité, de réduction ni de densité schématique.
\end{theoreme*}
\end{encadre}
{\small Les feuillets en esquissent une démonstration incomplète (page 3). Le lemme d'existence de $S_v$ y renvoie à \q{Invariants 1.2}. Mathlib n'a pas les schémas abéliens ni les groupes $p$-divisibles. On se reportera aux notes de Grothendieck (\lot{100}{1}{2}).\par}

\begin{encadre}
\begin{theoreme*}[Grothendieck, pages 22 à 24, descente des morphismes le long d'un épimorphisme propre\piste{100-proper-epimorphism-hom-descent}]\nonverifie
Soit $f : S' \to S$ un morphisme propre de schémas localement noethériens tel que $\mathcal O_S \to f_*\mathcal O_{S'}$ soit injectif, sans hypothèse de platitude, et soit $S'' = S' \times_S S'$. Pour tout $S$-schéma $X$ formellement non ramifié, la suite $\mathrm{Hom}_S(S, X) \to \mathrm{Hom}_S(S', X) \rightrightarrows \mathrm{Hom}_S(S'', X)$ est exacte ; le changement de base le long de $f$ est donc pleinement fidèle sur les schémas abéliens.
\end{theoreme*}
\end{encadre}
{\small Les feuillets en esquissent une démonstration incomplète (pages 22 et 24) : le lemme~1 se ramène, par l'algèbre $\mathrm{Ker}(f_*\mathcal O_{S'} \rightrightarrows h_*\mathcal O_{S''})$, au cas d'un épimorphisme fini et radiciel, le corollaire~1, dont la page~24 donne la réduction au cas local, l'image schématique et la finitude de $\mathcal O_{X,x}$ sur $\mathcal O_{S,s}$. La fin du corollaire~1 (non-ramification, Nakayama) et la radicialité de la réduction sont de l'édition, comme la représentabilité et la non-ramification du schéma des homomorphismes de schémas abéliens. Mathlib n'a pas la descente des morphismes le long d'un morphisme propre. On se reportera aux notes de Grothendieck (\lot{100}{2}{22}).\par}

\begin{encadre}
\begin{theoreme*}[Grothendieck, pages 30 à 38, algébrisation par les courbes\piste{100-formal-abelian-curve-criterion}]\nonverifie
Soit $R$ un anneau local noethérien complet. Un schéma abélien formel sur $\operatorname{Spf} R$ est algébrisable si et seulement si son changement de base à $\operatorname{Spf} B'$ l'est pour tout quotient intègre $B$ de dimension un de $R$, de normalisé $B'$, ou encore s'il le devient après un morphisme propre surjectif $S' \to \operatorname{Spec} R$.
\end{theoreme*}
\end{encadre}
{\small Les feuillets en esquissent une démonstration incomplète (pages 32 à 38). Mathlib n'a pas l'algébrisation des schémas formels. On se reportera aux notes de Grothendieck (\lot{100}{2}{30}).\par}

\begin{encadre}
\begin{theoreme*}[Pages 36 à 38, forme corrigée par l'édition, points de dimension un\piste{100-baire-dimension-one-points}]\nonverifie
Dans le spectre d'un anneau local noethérien complet, une famille dénombrable de fermés dont la réunion contient tous les points $x$ tels que $\dim \overline{\{x\}} \leq 1$ a une sous-famille finie qui recouvre le spectre.
\end{theoreme*}
\end{encadre}
{\small La page barre \q{complet}, que l'édition rétablit : sans cette hypothèse l'énoncé est faux. Les feuillets démontrent le cas de dimension $\leq 1$ et le corollaire, une propriété de Baire sur les $\mathfrak m^n$ ; le cas général de la proposition s'interrompt page~36 sur des lignes illisibles et annulées, et le passage du corollaire à la proposition n'est pas écrit. Le cas d'une variable du corollaire est le théorème~\ref{thm:100-points}.\par}

\begin{encadre}
\begin{theoreme*}[Grothendieck, pages 47 à 52, prolongement vers un modèle de Néron\piste{100-neron-extension-p-divisible}]\nonverifie
Sur un trait $S$, pour tout modèle de Néron $G$ et tout groupe $p$-divisible $M$ sur $S$, tout homomorphisme $M_K \to G_K$ des fibres génériques se prolonge de façon unique en $M \to G$, le théorème de pleine fidélité de Tate pour les groupes $p$-divisibles sur $S$ étant admis.
\end{theoreme*}
\end{encadre}
{\small La conclusion repose sur le théorème de Tate, et le cas où $G_K$ est extension d'une variété abélienne par un tore est traité dans le tapuscrit des pages 47 à~50, dont l'auteur n'est pas établi ; la part des feuillets manuscrits est le lemme des pages 51 et~52. Les feuillets en esquissent une démonstration incomplète : ce lemme renvoie deux fois à \q{on sait}, et la page~52 l'applique au modèle de Néron et conclut \q{On gagne !}.\par}
